% Lamina 10: escalon aditivo fijo (10/20/40 ppb), detector final y PARAAD-SA juntos, umbral p99.
% Datos: results/linea_base/rev24/metricas_banco.csv, escenario E1, nivel 0.01; filas
% Ens-Vb-X0.2596 (detector final) y PARAAD. Banco 2024-jul 2026, suma de 30 calendarios.
% Por cada tamano, cuatro barras: las dos de alarmas juntas (detector final azul, PARAAD-SA verde)
% y luego las dos de atacados. Alarmas = VP + FP (precision = VP / alto); atacados = VP + FN (recall).
% Mismo truco que estacion_atacada.tex: un solo eje "ybar stacked", cada barra con su x y las
% cinco series con las mismas 12 x (0 donde la pieza no va). Orden de apilado: VP del final,
% VP de PARAAD, FP (gris, comun), FN del final, FN de PARAAD.
% Eje Y en dias al año por estacion: "y filter" divide los conteos entre 630.27 estacion-años
% (23 005 estacion-dias / 365, por 10 calendarios por año); "axis cs" no pasa por el filtro.
% Generado con un script a partir del CSV; sustituye a estacion_atacada.tex y paraad_atacada_p99.tex.
\begin{tikzpicture}[
    etiqueta/.style={font=\sffamily\scriptsize, anchor=south, inner sep=1.5pt},
    par/.style={font=\sffamily\scriptsize\bfseries, text=texto, anchor=north},
    tamano/.style={font=\sffamily\small, text=texto, anchor=north},
  ]
\begin{axis}[
    ejes tesis,
    ybar stacked,
    bar width=0.2,                  % sin unidad = unidades del eje x
    scale only axis, width=12.0cm, height=3.6cm,
    xmin=0.42, xmax=3.98,
    y filter/.expression={y/630.27},  % estacion-dias -> dias al año por estacion
    ymin=0, ymax=14, ytick={0,2,...,12},
    ylabel={Días al año por estación},
    ylabel style={font=\sffamily\small},
    xtick={0.57,0.83,1.17,1.43,1.77,2.03,2.37,2.63,2.97,3.23,3.57,3.83},
    xticklabels={{\color{ensX}final},{\color{paraad}PARAAD},{\color{ensX}final},{\color{paraad}PARAAD},{\color{ensX}final},{\color{paraad}PARAAD},{\color{ensX}final},{\color{paraad}PARAAD},{\color{ensX}final},{\color{paraad}PARAAD},{\color{ensX}final},{\color{paraad}PARAAD}},
    xtick style={draw=none},
    xticklabel style={font=\sffamily\tiny\bfseries},
    legend style={at={(0.5,1.02)}, anchor=south, legend columns=4,
                  /tikz/every even column/.append style={column sep=0.45cm}},
    legend image code/.code={\fill[#1] (0cm,-0.1cm) rectangle (0.3cm,0.1cm);},
    clip=false,
  ]
  \addplot[barra=ensX] coordinates {(0.57,405) (0.83,0) (1.17,405) (1.43,0) (1.77,1477) (2.03,0) (2.37,1477) (2.63,0) (2.97,1703) (3.23,0) (3.57,1703) (3.83,0)};
  \addlegendentry{VP detector final}
  \addplot[barra=paraad] coordinates {(0.57,0) (0.83,180) (1.17,0) (1.43,180) (1.77,0) (2.03,1005) (2.37,0) (2.63,1005) (2.97,0) (3.23,1635) (3.57,0) (3.83,1635)};
  \addlegendentry{VP PARAAD-SA}
  \addplot[barra=sinataque!70!white] coordinates {(0.57,2998) (0.83,6033) (1.17,0) (1.43,0) (1.77,3198) (2.03,6042) (2.37,0) (2.63,0) (2.97,3471) (3.23,6071) (3.57,0) (3.83,0)};
  \addlegendentry{FP: falsa alarma}
  \addplot[barra=ensX!30!white] coordinates {(0.57,0) (0.83,0) (1.17,1372) (1.43,0) (1.77,0) (2.03,0) (2.37,300) (2.63,0) (2.97,0) (3.23,0) (3.57,74) (3.83,0)};
  \addlegendentry{FN: no detectado}
  \addplot[barra=paraad!30!white, forget plot] coordinates {(0.57,0) (0.83,0) (1.17,0) (1.43,1597) (1.77,0) (2.03,0) (2.37,0) (2.63,772) (2.97,0) (3.23,0) (3.57,0) (3.83,142)};

  % Precision sobre cada barra de alarmas y recall sobre cada barra de atacados (del CSV).
  \node[etiqueta, text=tenue] at (axis cs:0.57,3403/630.27) {P\,0.12};
  \node[etiqueta, text=tenue] at (axis cs:0.83,6213/630.27) {P\,0.03};
  \node[etiqueta, text=ensX] at (axis cs:1.17,1777/630.27) {R\,0.23};
  \node[etiqueta, text=paraad] at (axis cs:1.43,1777/630.27) {R\,0.10};
  \node[etiqueta, text=tenue] at (axis cs:1.77,4675/630.27) {P\,0.32};
  \node[etiqueta, text=tenue] at (axis cs:2.03,7047/630.27) {P\,0.14};
  \node[etiqueta, text=ensX] at (axis cs:2.37,1777/630.27) {R\,0.83};
  \node[etiqueta, text=paraad] at (axis cs:2.63,1777/630.27) {R\,0.57};
  \node[etiqueta, text=tenue] at (axis cs:2.97,5174/630.27) {P\,0.33};
  \node[etiqueta, text=tenue] at (axis cs:3.23,7706/630.27) {P\,0.21};
  \node[etiqueta, text=ensX] at (axis cs:3.57,1777/630.27) {R\,0.96};
  \node[etiqueta, text=paraad] at (axis cs:3.83,1777/630.27) {R\,0.92};

  % Separadores entre tamaños
  \draw[tenue!60, densely dashed, line width=0.4pt] (axis cs:1.6,0) -- (axis cs:1.6,13.5);
  \draw[tenue!60, densely dashed, line width=0.4pt] (axis cs:2.8,0) -- (axis cs:2.8,13.5);

  % Debajo: alarmas / atacados bajo cada par de barras y el tamaño del ataque
  \node[par] at ([yshift=-0.32cm]axis cs:0.700,0) {alarmas};
  \node[par] at ([yshift=-0.32cm]axis cs:1.300,0) {atacados};
  \node[tamano] at ([yshift=-0.72cm]axis cs:1.0,0) {10 ppb};
  \node[par] at ([yshift=-0.32cm]axis cs:1.900,0) {alarmas};
  \node[par] at ([yshift=-0.32cm]axis cs:2.500,0) {atacados};
  \node[tamano] at ([yshift=-0.72cm]axis cs:2.2,0) {20 ppb};
  \node[par] at ([yshift=-0.32cm]axis cs:3.100,0) {alarmas};
  \node[par] at ([yshift=-0.32cm]axis cs:3.700,0) {atacados};
  \node[tamano] at ([yshift=-0.72cm]axis cs:3.4,0) {40 ppb};
\end{axis}
\end{tikzpicture}
